\section{引言}

我们设计并实现了Google文件系统(GFS)，以满足Google数据处理需求
快速增长的压力。GFS与先前的分布式文件系统有许多共同目标，例如
性能、可扩展性、可靠性和可用性。然而，它的设计由我们对当前及预期
应用工作负载和技术环境的关键观察所驱动；这些观察与一些早期文件系统
设计假设存在显著差异。我们重新审视了传统选择，并在设计空间中探索了
截然不同的方向。

第一，组件故障是常态，而不是例外。该文件系统由数百乃至数千台采用
廉价通用部件构建的存储机器组成，并由数量相当的客户端机器访问。组件
的数量和质量几乎必然意味着：任意时刻总有部分组件无法工作，而且有些
组件无法从当前故障中恢复。我们见过由应用程序缺陷、操作系统缺陷、人为
错误，以及磁盘、内存、连接器、网络和电源故障造成的问题。因此，持续
监控、错误检测、容错和自动恢复必须成为系统不可或缺的一部分。

第二，按照传统标准来看，文件非常巨大。
多GB级别的文件很常见。每个文件通常包含许多应用层对象，比如网页文档。
当我们经常处理快速增长的、规模达到许多TB的数据集，并且这些
数据集由数十亿个对象组成时，即使文件系统能够支持数十亿个约KB大小的小文件，
管理这些小文件也会非常棘手。因此，像I/O操作大小、块大小这样的设计假设和参数，
都必须重新审视。

第三，大多数文件通过追加新数据而不是覆盖已有数据来修改。文件内部的
随机写入实际上几乎不存在。文件一旦写入，之后只会被读取，而且通常只会
被顺序读取。多种数据都具有这些特征：有些数据可能形成大型数据仓库，
供数据分析程序进行顺序扫描处理；有些可能是运行中应用持续生成的数据流；
有些可能是归档数据；还有些可能是在一台机器上产生、由另一台机器同时
或稍后处理的中间结果。鉴于超大文件的这种访问模式，追加成为性能优化和
原子性保证的重点，而在客户端缓存数据块的吸引力则降低了\emph{(译者
认为这里说的是传统客户端会尝试缓存数据，但是对于GFS这类通常顺序
读取大量连续数据块的程序，缓存很不占优势：首先是文件太大，其次是
缺少时间局部性，内容多半顺序读取一次后续就不读了，
所以客户端完全不缓存数据块内容本身)}。

第四，对应用程序和文件系统API进行协同设计，能够通过提高灵活性使整个
系统受益。例如，我们放宽了GFS的一致性模型，在不向应用程序施加沉重负担
的前提下，大幅简化了文件系统。我们还引入了原子追加操作，使多个客户端
无需彼此额外同步，就能并发向同一文件追加数据。本文后面会更详细地讨论
这些内容。

目前，多个GFS集群已针对不同用途部署。最大的集群拥有超过1000个存储
节点、超过300TB的磁盘存储，并持续受到数百个不同机器上的客户端的高频
访问。